\clearpage\pagestyle{otMain}\setlength{\otBlocEspacement}{2.5mm}\fontsize{11}{15}\selectfont\corrigesChapitre{19}{Hasard et probabilités}
\begin{blocLivre}{cor-19-1}\begin{corrige}{1}{Les issues d’un dé}\demarche{Préciser ce qu’on observe.}
Les issues sont $\res{1,2,3,4,5,6}$. Une issue décrit ici le numéro de la face obtenue.
\end{corrige}\end{blocLivre}
\begin{blocLivre}{cor-19-2}\begin{corrige}{2}{Une pièce équilibrée}\demarche{Identifier les issues équiprobables.}
Les deux faces sont équiprobables : $\res{1/2=0{,}5=50\,\%}$.
\end{corrige}\end{blocLivre}
\begin{blocLivre}{cor-19-3}\begin{corrige}{3}{Une échelle de probabilité}\demarche{Comparer l’évènement à toutes les issues.}
8 n’est pas une issue : A est impossible, probabilité $\res0$. Toutes les faces sont inférieures à 7 : B est certain, probabilité $\res1$. C a pour probabilité $\res{1/2}$. D est très probable, très près de 1 sans être certain ; E est très peu probable, très près de 0 sans être impossible.
\label{fig-cor-19-3}\figureGeo[corrige]{proba-3}
\end{corrige}\end{blocLivre}
\begin{blocLivre}{cor-19-4}\begin{corrige}{4}{Une face donnée}\demarche{Compter les issues favorables.}
Une seule issue sur six convient : $\res{1/6}$. Cela ne signifie pas qu’un 5 apparaîtra exactement une fois tous les six lancers.
\end{corrige}\end{blocLivre}
\begin{blocLivre}{cor-19-5}\begin{corrige}{5}{Lire une chance sur quatre}\demarche{Convertir la fraction en pourcentage.}
$1/4=25/100=\res{25\,\%}$. Le nombre quatre est le total des parts équiprobables.
\end{corrige}\end{blocLivre}
\begin{blocLivre}{cor-19-6}\begin{corrige}{6}{Un évènement à plusieurs issues}\demarche{Lister avant de compter.}
Les issues favorables sont 2, 4 et 6 : $3/6=\res{1/2}$. Trois issues différentes peuvent réaliser le même évènement.
\end{corrige}\end{blocLivre}
\begin{blocLivre}{cor-19-7}\begin{corrige}{7}{Une urne}\demarche{Compter les objets, pas seulement les couleurs.}
Il y a $4+3+1=8$ boules, dont trois bleues. La probabilité vaut $\res{3/8=37{,}5\,\%}$. Les boules sont équiprobables, pas les couleurs.
\label{fig-cor-19-7}\figureGeo[corrige]{proba-7}
\end{corrige}\end{blocLivre}
\begin{blocLivre}{cor-19-8}\begin{corrige}{8}{Une roue régulière}\demarche{Utiliser l’égalité des secteurs.}
Les secteurs 6, 7 et 8 conviennent, sur huit secteurs équiprobables : $\res{3/8}$.
\label{fig-cor-19-8}\figureGeo[corrige]{proba-8}
\end{corrige}\end{blocLivre}
\begin{blocLivre}{cor-19-9}\begin{corrige}{9}{Des jetons identifiés}\demarche{Définir précisément l’évènement.}
Les jetons R1, B1 et V1 conviennent : $3/6=\res{1/2}$. Leur couleur ne change pas le nombre inscrit.
\label{fig-cor-19-9}\figureGeo[corrige]{proba-9}
\end{corrige}\end{blocLivre}
\begin{blocLivre}{cor-19-10}\begin{corrige}{10}{Fréquence observée}\demarche{Séparer observation et probabilité supposée.}
La fréquence vaut $18/40=\res{0{,}45=45\,\%}$. Cet écart à 50 \% ne suffit pas à prouver un déséquilibre ; les résultats d’une série aléatoire peuvent varier.
\end{corrige}\end{blocLivre}
\begin{blocLivre}{cor-19-11}\begin{corrige}{11}{Préparer une expérience}\demarche{Conserver les mêmes conditions d’essai.}
Tirer sans regarder, noter la couleur, remettre la boule puis mélanger, et recommencer. La remise maintient deux rouges sur cinq à chaque essai : la probabilité du modèle reste $\res{2/5}$. Sans remise, la composition change.
\label{fig-cor-19-11}\figureGeo[corrige]{proba-11}
\end{corrige}\end{blocLivre}
\begin{blocLivre}{cor-19-12}\begin{corrige}{12}{Comparer deux jeux}\demarche{Comparer deux fractions.}
A donne $1/6=5/30$ ; B donne $3/10=9/30$. La probabilité est plus grande dans le jeu $\res B$. Cela ne garantit pas le gain au prochain essai.
\end{corrige}\end{blocLivre}
\begin{blocLivre}{cor-19-13}\begin{corrige}{13}{Construire une urne}\demarche{Passer de la proportion à l’effectif.}
Un tiers de douze vaut $\res4$ : il faut quatre rouges et $\res8$ bleues. Contrôle : $4/12=1/3$.
\label{fig-cor-19-13}\figureGeo[corrige]{proba-13}
\end{corrige}\end{blocLivre}
\begin{blocLivre}{cor-19-14}\begin{corrige}{14}{Une fausse compensation}\demarche{Distinguer tendance de fréquence et résultat imposé.}
Non. Si les conditions restent celles d’une pièce équilibrée à chaque lancer, pile et face ont encore chacun une probabilité $\res{1/2}$. Une série passée n’impose pas une compensation au lancer suivant.
\end{corrige}\end{blocLivre}
\begin{blocLivre}{cor-19-15}\begin{corrige}{15}{Deux relevés différents}\demarche{Calculer les écarts sans transformer une tendance en garantie.}
Les fréquences sont 0,6 et 0,46 ; leurs écarts à 0,5 valent $\res{0{,}10}$ et $\res{0{,}04}$. La seconde est ici plus proche, mais chaque série plus longue n’est pas obligée de l’être : les fluctuations restent possibles.
\end{corrige}\end{blocLivre}
